% !TeX root = ../main.tex
\section{结果与局限}\label{sec:results}
图~\ref{fig:comparison} 的子图~\ref{fig:baseline} 和~\ref{fig:calibrated}
分别展示两种方法；表~\ref{tab:metrics} 给出汇总。
Calibrated 的 RMSE 为 $0.500\,{}^\circ\mathrm{C}$，
比 Baseline 的 $1.000\,{}^\circ\mathrm{C}$ 低 50\%。
这一比较仅对当前六条合成样本成立，不代表真实环境中的泛化效果。
\begin{figure}[htbp]
 \centering
 \begin{subfigure}{0.48\linewidth}
  \centering\includegraphics[width=\linewidth]{baseline.pdf}
  \caption{Baseline}\label{fig:baseline}
 \end{subfigure}\hfill
 \begin{subfigure}{0.48\linewidth}
  \centering\includegraphics[width=\linewidth]{calibrated.pdf}
  \caption{Calibrated}\label{fig:calibrated}
 \end{subfigure}
 \caption{两种方法的误差对比（教学合成数据）。}\label{fig:comparison}
\end{figure}
\begin{table}[htbp]
 \centering
 \caption{温度误差汇总（教学合成数据）。}\label{tab:metrics}
 \begin{tabular}{lrr}
  \toprule
  \multicolumn{1}{c}{方法} & \multicolumn{2}{c}{误差（$^\circ$C）}\\
  \cmidrule(lr){2-3}
   & RMSE & MAE\\
  \midrule
  Baseline & 1.000 & 1.000\\
  Calibrated & 0.500 & 0.500\\
  \bottomrule
 \end{tabular}
\end{table}
